\documentclass[12pt]{article}

\usepackage{sbc-template}
\usepackage{graphicx,url}
\usepackage[utf8]{inputenc}
\usepackage[english]{babel}
\usepackage{booktabs}

\sloppy

\title{Toward the Evolution of the Desktop GUI:\\ Human-Centered AI as a Paradigm for\\ Rethinking Operating System Interaction}

\author{Jose Ramon Aragon Toledo\inst{1}, Ricardo Ruíz Rodríguez\inst{1}}

\address{Universidad Tecnológica de la Mixteca (UTM)\\
  Huajuapan de León, Oaxaca -- México
  \email{\{jaragon,rruiz\}@mixteco.utm.mx}
}

\begin{document}

\maketitle

\begin{abstract}
Innovation in desktop operating system interfaces has stagnated over the past two
decades. Despite exponential growth in computational power, tools such as file
explorers, terminals, and configuration managers have scarcely evolved since the
1990s. This plateau particularly affects GNU/Linux, whose power remains inaccessible
to non-expert users due to command-line interaction models from the 1970s. This paper
analyzes desktop interaction stagnation through the lens of Human-Centered AI (HCAI)
and presents BOTAS, a voice-based assistant enabling Spanish-speaking users to
interact with GNU/Linux via natural language. Evaluation across 200 tasks shows
89\% command-generation accuracy and 100\% detection of destructive operations.
\end{abstract}


\section{Introduction}
\label{sec:intro}

We operate with processors a hundred million times faster than those of the 1980s,
yet the fundamental interaction with a desktop operating system has barely changed
in twenty years \cite{Jenson2024}. Scott Jenson---a veteran UX designer who worked
on the original Macintosh Finder at Apple and later at Google---argues that ``this
immobility is not a technical limitation but a capitulation of design''
\cite{Jenson2024}. The industry has moved from the creative transformation that
defined early computing---where Xerox PARC concepts were radically reimagined---to
shallow cross-platform imitation.

GNU/Linux powers approximately 90\% of cloud infrastructure and 100\% of the
world's top 500 supercomputers \cite{W3Techs2023, Top5002023}, yet its command-line
interface perpetuates a 1970s paradigm \cite{Raymond2003} demanding cryptic syntax
with no feedback about operational risk. Meanwhile, commercial voice
assistants handle trivial queries but none addresses deep OS interaction
\cite{Luger2016}, and tools like ChatGPT suggest commands without executing them
or validating safety. This paper argues that Human-Centered AI (HCAI)
\cite{Shneiderman2020, Shneiderman2022} provides a framework for overcoming this
stagnation, and presents BOTAS, a voice-based assistant that materializes HCAI
principles for GNU/Linux interaction in Spanish.

This work contributes: (1)~an analysis of desktop interface stagnation through the
HCAI lens; (2)~the presentation of BOTAS, whose evolution from minimum viable
prototype to modular system illustrates HCAI in practice; and (3)~preliminary
evaluation with accuracy metrics, safety detection, and pilot usability feedback.


\section{Conceptual Framework and Related Work}
\label{sec:framework}

\subsection{Stagnation of Desktop Interfaces}
\label{subsec:stagnation}

HCI has traversed periods of genuine innovation---the desktop metaphor, the Web,
the touch revolution---that redefined how users form mental models of computing
\cite{Norman1988}. On the desktop, however, the past two decades constitute a
functional plateau. Jenson \cite{Jenson2024} illustrates this with an anecdote:
in the early 1980s, while working on the Finder, his team noticed that truncated
filenames lost critical information with end-placed ellipses. ``I suggested placing
the ellipsis in the middle to preserve the file extension---it was implemented in
two days and, forty years later, that code is still running'' \cite{Jenson2024}.
Today, in the era of large language models, operating systems still employ
file-navigation paradigms from three decades ago---a paradox especially acute for
GNU/Linux, where the CLI constitutes the primary barrier to a computing platform
that dominates professional infrastructure \cite{W3Techs2023}.

\subsection{Human-Centered AI}
\label{subsec:hcai}

Shneiderman \cite{Shneiderman2020, Shneiderman2022} formalized HCAI as a paradigm
where AI systems \textit{amplify} human capabilities rather than replace them,
built on four principles: \textbf{transparency} (users understand system actions),
\textbf{human control} (users supervise and can reject proposals),
\textbf{safety} (prevention of harmful actions, even if inadvertently
requested---critical where a typo can cause irreversible data loss), and
\textbf{empowerment} (users progressively learn underlying mechanisms, building
autonomy over dependency). Amershi et al. \cite{Amershi2019} complement this with
18 practical guidelines for human-AI interaction. As AI adoption grows and users
discover the possibilities these tools offer, HCAI provides a coherent
foundation for rethinking OS interaction: an assistant grounded in these
principles can make GNU/Linux accessible while ensuring users learn rather than
depend.

\subsection{Voice Assistants and the OS Interaction Gap}
\label{subsec:voice-gap}

Advances in speech recognition---particularly Whisper \cite{Radford2022}---and
in LLMs \cite{Brown2020, OpenAI2023} have reached practical multilingual maturity.
NL2Bash \cite{Lin2018} explored natural language to Bash translation (36\%
exact-match accuracy) but without real-time interaction or safety validation. No
existing system integrates voice recognition, command execution, safety
classification, and educational feedback within an HCAI workflow. This
integration gap defines the opportunity this work addresses.


\section{Proposal: BOTAS with HCAI Approach}
\label{sec:proposal}

HCAI principles are materialized on GNU/Linux because this is where the
intervention yields the greatest impact: a powerful system whose main adoption
barrier is interaction, not functionality \cite{Warschauer2004}. Its openness
permits deep integration, and its zero cost eliminates economic barriers---decisive
for educational contexts in Latin America.

\subsection{Origins and Evolution}
\label{subsec:origins}

BOTAS (\textit{Bot de Operaciones y Tareas Automatizadas del Sistema}) originated
during a social service program focused on AI-integrated applications at a Mexican
public university. Initially a minimum viable product, it explored whether natural
language could bridge novice users and the command line. Early versions relied
entirely on cloud services (Whisper for transcription, GPT-4o for interpretation)
with a GTK3 interface and local TTS via eSpeak. Through iterative development,
the system incorporated: offline wake-word detection (Vosk), a local
intent-recognition module resolving routine requests without LLM calls, automatic
distribution detection across six Linux families, a visual state indicator, and
structured request logging. The strategy of routing all requests through the LLM
was supplemented by local detection, reducing latency and costs; rigid error
handling is being redesigned toward clarification dialogues.

\subsection{Current Capabilities}
\label{subsec:capabilities}

BOTAS implements a five-layer pipeline: (1)~input via GUI or wake-word;
(2)~local intent detection through regex patterns; (3)~GPT-4o command
parsing for complex requests; (4)~deterministic safety validation at three risk
levels (\textit{safe}, \textit{warning}, \textit{dangerous}); and (5)~execution
with educational feedback and TTS narration. Supported operations include file
management (creation, renaming, moving, search, deletion, editing), system
queries (memory, disk, processes, ports, time), process control, networking,
permissions, and web search. Every operation includes an explanation of the
generated command, fulfilling the empowerment principle. The safety layer requires
explicit confirmation for risky operations. Distribution detection (Debian,
Fedora, Arch, openSUSE, Alpine, Void) translates package-management commands
automatically.

\subsection{Scope and Ongoing Developments}
\label{subsec:scope}

The system currently depends on cloud services (Whisper, GPT-4o). We are
developing local operation using DeepSeek-R1-0528-Distill-8b---served via local
network or on the host machine---with Vosk for speech recognition, eliminating
cloud dependency and addressing privacy concerns for institutional environments.
An expected accuracy decrease relative to cloud models is acknowledged.

Recognized limitations include difficulty with ambiguous requests, multi-command
concatenation in single utterances, technical jargon disambiguation, and absent
clarification dialogues. Features under development include improved error
recovery, responsive interface layout, and cross-distribution installation
tooling. We plan to release BOTAS as open-source software; post-release feedback
will inform version~2.0.


\section{Preliminary Evaluation}
\label{sec:evaluation}

\subsection{Functional Testing}
\label{subsec:testing}

BOTAS was evaluated on 200 representative Linux tasks across five categories,
specified as Spanish natural-language requests. The dataset was compiled and
verified manually by the developer---a constraint we acknowledge, as it may
introduce confirmation bias. Table~\ref{tab:accuracy} summarizes results.

\begin{table}[h]
\caption{Command generation accuracy by task category}
\label{tab:accuracy}
\centering
\begin{tabular}{lcc}
\toprule
\textbf{Category} & \textbf{Tasks} & \textbf{Accuracy} \\
\midrule
File management     & 45 & 84\% \\
System information  & 40 & 85\% \\
Process control     & 35 & 83\% \\
Networking          & 30 & 93\% \\
Permissions         & 25 & 80\% \\
\midrule
\textbf{Total}      & \textbf{200} & \textbf{89\%} \\
\bottomrule
\end{tabular}
\end{table}

Error modes: ambiguous phrasing (6\%), multi-pipe flag errors (3\%),
distribution-specific unavailability (2\%). Errors were categorized as
transcription failures, LLM misinterpretation, or execution-state mismatches.
The deterministic safety layer detected 100\% of 43 dangerous operations
(recursive deletion, privilege escalation, system modifications, permission
changes)---a rate guaranteed by regex-based pattern matching independent of model
variability. The local intent detector resolved 29\% of tasks without LLM calls,
at 100\% accuracy, saving 500--1500\,ms per request.

\subsection{Pilot Usability Study}
\label{subsec:usability}

Three users with no prior GNU/Linux experience evaluated BOTAS at its MVP stage.
While insufficient for statistical generalization---acknowledged
prominently---this pilot yielded qualitative insights. Participants found the
natural-language execution capability, educational explanations, and safety
feedback encouraging; these positive impressions from users who had never used a
Linux terminal motivated continued development and this contribution to the HCI
community. Constructive feedback identified confusion with verbal syntax,
frustration with seemingly excessive safety confirmations, and difficulty with
error messages---directly informing the redesign toward clarification dialogues.
Latency data: Whisper transcription averages 1--3\,s; Vosk wake-word detection
responds in under 200\,ms.


\section{Discussion and Conclusions}
\label{sec:discussion}

Desktop interface stagnation is a design problem, not a technical one: the
enabling technologies exist and are mature. HCAI provides a theoretical and
practical framework for addressing it, prioritizing user autonomy over opaque
automation. From a Latin American perspective, the linguistic barrier excludes
users whose documentation is in English, and closed-ecosystem costs burden
educational institutions \cite{Warschauer2004}. Voice interaction in Spanish
reduces these barriers without requiring prior CLI expertise, offering guided
learning rather than conditioning through repetition.

Future work includes consolidating local operation (DeepSeek-R1 with Vosk),
a formal usability study with an expanded sample, cross-distribution validation,
and the open-source release.

The HCAI principles of transparency, control, safety, and
empowerment---materialized in BOTAS---demonstrate that access barriers to
powerful systems can be reduced without sacrificing capabilities. The goal is
not to make technology simpler, but to make it more accessible---and in that
process, more human.


\bibliographystyle{sbc}
\bibliography{paper-en}

\end{document}
